% theme: dot_product_angles_and_perpendicularity
\MajorQuestionLesson{内積となす角・垂直}
\SetRowSpace{5mm}

\TypeQuestionSingle{次の2つのベクトルのなす角 $\theta$ を求めなさい。}{angle_between_vectors}
\QQRow{$\begin{pmatrix}1\\0\end{pmatrix},\ \begin{pmatrix}1\\\sqrt3\end{pmatrix}$}{}{$\begin{pmatrix}1\\1\end{pmatrix},\ \begin{pmatrix}-1\\1\end{pmatrix}$}{}
\QQRow{$\begin{pmatrix}3\\4\end{pmatrix},\ \begin{pmatrix}4\\-3\end{pmatrix}$}{}{$\begin{pmatrix}2\\-2\end{pmatrix},\ \begin{pmatrix}-1\\0\end{pmatrix}$}{}
\QQRow{$\begin{pmatrix}\sqrt3\\1\end{pmatrix},\ \begin{pmatrix}\sqrt3\\-1\end{pmatrix}$}{}{$\begin{pmatrix}-2\\2\sqrt3\end{pmatrix},\ \begin{pmatrix}1\\\sqrt3\end{pmatrix}$}{}

\TypeQuestionDouble{次の2つのベクトルが垂直になるように、実数 $k$ の値を求めなさい。}{perpendicular_parameter}
\QQRow{$\begin{pmatrix}2\\3\end{pmatrix},\ \begin{pmatrix}k\\-2\end{pmatrix}$}{}{$\begin{pmatrix}k\\4\end{pmatrix},\ \begin{pmatrix}2\\-3\end{pmatrix}$}{}
\QQRow{$\begin{pmatrix}k-1\\2\end{pmatrix},\ \begin{pmatrix}3\\k\end{pmatrix}$}{}{$\begin{pmatrix}1\\k+2\end{pmatrix},\ \begin{pmatrix}2k\\-1\end{pmatrix}$}{}

\TypeQuestionSingle{次のベクトルに垂直な単位ベクトルをすべて求めなさい。}{perpendicular_unit_vector}
\QQRowThree{$\begin{pmatrix}3\\4\end{pmatrix}$}{}{$\begin{pmatrix}-5\\12\end{pmatrix}$}{}{$\begin{pmatrix}1\\-1\end{pmatrix}$}{}
\QQRowThree{$\begin{pmatrix}0\\7\end{pmatrix}$}{}{$\begin{pmatrix}2\\\sqrt5\end{pmatrix}$}{}{$\begin{pmatrix}-8\\-6\end{pmatrix}$}{}

\LessonPageBreak
\SetRowSpace{5mm}

\TypeQuestionDouble{$|\vec{a}|,|\vec{b}|$ と $\vec{a}\cdot\vec{b}$ が次のように与えられている。2つのベクトルのなす角を求めなさい。}{angle_from_dot_product}
\QQRow{$|\vec{a}|=2,\ |\vec{b}|=3,\ \vec{a}\cdot\vec{b}=3$}{}{$|\vec{a}|=4,\ |\vec{b}|=5,\ \vec{a}\cdot\vec{b}=-10$}{}
\QQRow{$|\vec{a}|=3,\ |\vec{b}|=6,\ \vec{a}\cdot\vec{b}=0$}{}{$|\vec{a}|=\sqrt2,\ |\vec{b}|=\sqrt6,\ \vec{a}\cdot\vec{b}=\sqrt3$}{}

\TypeQuestionSingle{次の条件を満たす実数 $k$ を求めなさい。}{specified_angle_parameter}
\QQRow{$\begin{pmatrix}1\\0\end{pmatrix}$ と $\begin{pmatrix}k\\1\end{pmatrix}$ のなす角が $45^\circ$}{}{$\begin{pmatrix}1\\\sqrt3\end{pmatrix}$ と $\begin{pmatrix}k\\1\end{pmatrix}$ のなす角が $60^\circ$}{}
\QQRow{$\begin{pmatrix}2\\1\end{pmatrix}$ と $\begin{pmatrix}1\\k\end{pmatrix}$ のなす角が $90^\circ$}{}{$\begin{pmatrix}1\\-1\end{pmatrix}$ と $\begin{pmatrix}k\\2\end{pmatrix}$ のなす角が $135^\circ$}{}

\LessonPageBreak
\SetRowSpace{6mm}

\TypeQuestionDouble{次の2つのベクトルのなす角が鋭角、直角、鈍角のいずれであるかを答えなさい。}{classify_angle_by_dot_product}
\QQRow{$\begin{pmatrix}2\\1\end{pmatrix},\ \begin{pmatrix}-1\\3\end{pmatrix}$}{}{$\begin{pmatrix}3\\-2\end{pmatrix},\ \begin{pmatrix}2\\3\end{pmatrix}$}{}
\QQRow{$\begin{pmatrix}-4\\1\end{pmatrix},\ \begin{pmatrix}1\\2\end{pmatrix}$}{}{$\begin{pmatrix}-2\\-3\end{pmatrix},\ \begin{pmatrix}-4\\1\end{pmatrix}$}{}


\TypeQuestionDouble{次の3点について、指定された角を求めなさい。}{angle_from_three_points}
\QQ{$\mathrm{A}(3,-2),\ \mathrm{B}(1,-2),\ \mathrm{C}(3,0)$ の $\angle\mathrm{ABC}$}{}
\QQ{$\mathrm{A}(0,3),\ \mathrm{B}(-1,2),\ \mathrm{C}(-2,3)$ の $\angle\mathrm{ABC}$}{}
\QQ{$\mathrm{A}(3,1),\ \mathrm{B}(2,1),\ \mathrm{C}(0,3)$ の $\angle\mathrm{ABC}$}{}
\QQ{$\mathrm{A}(-3,-2),\ \mathrm{B}(-2,-1),\ \mathrm{C}(0,1)$ の $\angle\mathrm{ABC}$}{}

\TypeQuestionDouble{$\vec{a},\vec{b}$ は零ベクトルでないとする。次の条件から、$\vec{a}$ と $\vec{b}$ が垂直であるかどうかを判定しなさい。}{perpendicular_from_norms}
\QQRow{$|\vec{a}+\vec{b}|=|\vec{a}-\vec{b}|$}{}{$|\vec{a}+\vec{b}|^2=|\vec{a}|^2+|\vec{b}|^2$}{}
\QQRow{$|\vec{a}+\vec{b}|=|\vec{a}|+|\vec{b}|$}{}{$|\vec{a}-\vec{b}|^2=|\vec{a}|^2+|\vec{b}|^2$}{}
